% Fyll ut {} øverst. Tomme felt viser hjelpetekst. Bruk LuaLaTeX eller pdfLaTeX.
\documentclass[10pt,a4paper,landscape]{article}

\ExplSyntaxOn
\prop_new:N \g_plan_prop
\NewDocumentCommand{\fyllinn}{m +m}
  {\prop_gput:Nnn \g_plan_prop {#1} {#2}}
\ExplSyntaxOff
% Grunnopplysninger
\fyllinn{Dato}{19.10.26}
\fyllinn{Trinn}{VG1}
\fyllinn{Fag}{1T}
\fyllinn{Varighet}{2 timer}
% Kompetansemål
% Ett mål per \item.
\fyllinn{Kompetansemål}{%
  \item forklare forskjellen mellom en identitet, en ligning,
        et algebraisk uttrykk og en funksjon
  \item utforske strategier for å løse ligninger, ligningssystemer og
        ulikheter og argumentere for tenkemåtene sine
  \item lese og forstå matematiske bevis og utforske og utvikle bevis
        i relevante matematiske emner
}
% Mål og innhold
% Hva skal elevene kunne, forstå eller mestre etter undervisningen?
\fyllinn{Læringsmål}{ 

    Eleven skal bli kjent med forkorting av uttrykk av høyere grad. 
    Eleven skal kunne forstå hvordan faktorisering og produktregelen kan brukes for å løse ligninger.  

}
% Hvilke begreper, temaer og ferdigheter skal undervisningen handle om?
\fyllinn{Faglig innhold}{ Bruke og forstå kvadratsetningene $(a+ b)^2 = a^2 + 2ab + b^2$. Bli introdusert til egenskapene for andregradsuttrykk i ligninger. Elevene skal også bli kjent med begrepet faktorisering av uttrykk og hvordan dette henger
sammen med produktregelen der $ a\cdot b = 0 \implies a = 0 \textit{ eller } b=0 $. Eleven skal også kunne faktorisere enkle andregradsuttrykk. }
% Hvorfor har du valgt disse læringsmålene og dette faglige innholdet?
\fyllinn{Begrunnelse for mål}{ Elevene har brukt mange undervisningstimer på å utlede første og andre kvadratsetning. De har nå god trening i å gå fra faktorisert form til et utvidet uttrykk. I denne dobbeltimen skal vi prøve å lære elevene hvordan de kan gå fra et utvidet uttrykk til faktorisert form. For å motivere elevene bruker vi produktregelen og ligninger som utgangspunkt, slik at elevene må faktorisere uttrykket for å finne svaret. En av fallgruvene er at produktregelen, som algebraisk sett er et enkelt resultat, kan virke svært abstrakt når vi kombinerer den med faktorisering av andregradsuttrykk.
}
% Elevgruppen og vurdering
% Hvilke forkunnskaper har elevene? Hvilke behov for tilpasning ser du?
\fyllinn{Læreforutsetninger}{ I denne klassen har elevene variert bakgrunn, da de kommer fra ulike ungdomskoler, sosiale livssituasjoner, ulike lærere og ulik praksis. Vi kan dog forvente at de er kjent med ligninger og er kjent med
regneoperasjoner, men det vil antageligvis kreve oppfrisking og endring av misforståelser. Det kan også være variasjon i disiplin og utholdenhet, da mange kan ha et svært kort oppmerksomhetsvindu når ting blir vanskelig. Det vil si at det kan
eksistere svakere elever med både lav utholdenhet og som sliter med grunnleggende regneferdigheter i rommet.   }
% Hvorfor passer tilpasningene til denne elevgruppen?
\fyllinn{Begrunnelse for tilpasning}{ 

     % Det er verdt å nevne at dette er et fag som viser store krav til tilegnelse av abstrakt informasjon, og elevene har tatt et aktivt valg å velge dette grunnet utdanningsmuligheter, nyskjerrighet og ambisiøsitet. 

    For å bygge god basis ønsker vi å utfordre elevene til å bygge generell utholdenhet i å tilegne seg ny abstrakt informasjon, men samtidig ikke presse de svake elevene for mye slik at de oppnår et stadium av hjelpesløshet for tidlig i undervisningen.
    Derfor er undervisningen bygget opp ved å gradvis bygge klosser med logiske premisser. De sterke
    elevene bygger også opp forståelsen da de kan utforske nyttigheten av produktregelen og dens begrensninger.        
}
% Hvordan får du informasjon om elevenes læring? Hva skal du se etter?
% Hva skal eventuelt vurderes, og hvordan gir du tilbakemeldinger?
\fyllinn{Vurdering}{ Det blir også stilt muntlige kontrollspørsmål kontinuerlig gjennom tavleundervisningen, både til svake og sterke elever. Vi sjekket også utregningene på de gule lappene for å fange opp vanlige feil. På vertikal tavle var
vurderingen muntlige forklaringer mellom elevene.}
% Hvorfor velger du disse vurderingsformene?
\fyllinn{Begrunnelse for vurdering}{ Effektiv og sosial måte. Da dette var en undervisning der vi lærte om nye verktøy i matematikk, det vil si å kombinere faktorisering og produktregelen, forventet vi ikke at elevene skulle forstå det i denne timen.
Et av målene var å teste elevenes utholdenhet idet de tilegner seg abstrakt informasjon.       } % Gjennomføring
% Når skjer de ulike aktivitetene, og hvor lenge varer de?
\fyllinn{Tidsplan}{%
  \textbf{1. time}\par\smallskip
  10 min\par
  1a)--1b)\par\smallskip
  10 min\par
  1c)\par\smallskip
  25 min\par
  1d)\par\medskip
  \textbf{Pause}\par\medskip
  \textbf{2. time}\par\smallskip
  10 min\par
  2a)--2b)\par\smallskip
  10 min\par
  2c)\par\smallskip
  25 min\par
  Oppgaver\par
  i boka
}
% Hva gjør læreren? Hvordan organiseres undervisningen?
\fyllinn{Lærerens aktiviteter}{

Første time var lette eksempler for å illustrere hva vi skal oppnå. 
I dette tilfellet var progresjonen å forenkle

\[
\begin{aligned}
    \text{1a)}\quad &\frac{4x}{x}  ,
    \qquad \text{1b)}\quad \frac{x(x+1)}{(x+1)} , \\
    \text{1c)}\quad &\frac{x^2 + x}{x+1}  \\
\end{aligned}
\]



Elevene ble da satt til å løse oppgaven på såkalt vertikal tavle, to og to:
\begin{equation*}
    \text{1d) }\frac{x^2 + 2x + 1}{x+1}
\end{equation*}

 I andre time gikk vi frem på samme måte for å introdusere produktregelen, med følgende progresjon:

\[
\begin{aligned}
    \text{2a)}\quad & 4x = 0,
    \qquad \text{2b)}\quad 4(x-1) = 0, \\
    \text{2c)}\quad &(x-1)(x-1)= 0
\end{aligned}
\]

Etterpå ble de satt til å løse oppgaver i boka der de skulle jobbe alene med tilsvarende oppgaver. 
}

% Hva gjør elevene? ArbeiJder de alene, i grupper eller i fellesskap?
\fyllinn{Elevenes aktiviteter}{I oppgave 1a) og 1b) gikk vi for tradisjonell tavleundervisning der elevene følger med på tavla. I oppgave 1c) fikk alle elevene en gul lapp der de selvstendig skulle finne et svar, før vi gikk gjennom fasiten på tavla. 

    
}
% Hvilke læremidler og ressurser brukes? F.eks. lærebok og digitale verktøy.
\fyllinn{Læringsressurser}{ Tavleundervisning, vertikal tavle, gul lapp}
% Hvorfor velger du disse arbeidsformene, aktivitetene og ressursene?
\fyllinn{Begrunnelse for aktiviteter}{

Allerede i oppgave $1c)$ mistenkte vi at elevene måtte tilegne seg ny kunnskap, da de må kunne faktorisere $x^2 + x$, og derfor
    valgte vi å gi elevene en gul lapp hver der de skulle prøve å gjøre den spesifikke oppgaven i stillhet. Denne gule lappen er i prinsippet et verktøy for at de her og nå skal forsøke å løse oppgaven med egne ord og med egne tanker. 


    Vertikal tavle i dette tilfellet er for å står sammen om å bruke kreativitet og og problemløsning seg gjennom nye oppgaver.  

    De hadde allerede i første time lært nye nytt, og i mange tilfeller ville en lærer satt stopp og repetere hva vi har lært. Da vi vil teste deres utholdenhet valge vi da å introdusere produktregelen og faktorisering, som er nye abstrakt verktøy
    for mange. Dette er viktig for å sette dem i en posisjon der de er tvungen til å tenke, selv etter å ha lært mye allerede fra første time. Senere ble dem sendt å løse oppgaver i boka inviduelt. 

}
% Veiledning
% Hva ønsker du veiledning på? Hvilke spørsmål vil du ta opp?
\fyllinn{Spørsmål til veiledning}{
    I andre time gikk tråkket vi bevisst utenfor elevenes komfortsone, i en situasjon der mange i klassen ikke er i stand til å løse problemet. Kan vi potensielt sette for mye press på svakere elever? 

}
% Felt for veilederens tilbakemeldinger.
\fyllinn{Veileders kommentarer}{
    Ja, (hjelpesløshet, utmattelse, mestringsglede, potensielt ueffektiv læringmilkø pga høyt tempo.) 

}
% OPPSETT -- ingen utfylling nedenfor.

\usepackage{iftex}
\ifPDFTeX
  \usepackage[utf8]{inputenc}
  \usepackage[T1]{fontenc}
  \IfFileExists{phvr8t.tfm}{\usepackage[scaled=0.95]{helvet}}{}
\else
  \usepackage{fontspec}
  \IfFontExistsTF{Arial}{\setsansfont{Arial}}{%
    \IfFontExistsTF{TeX Gyre Heros}{\setsansfont{TeX Gyre Heros}}{%
      \setsansfont{Nimbus Sans}}}
\fi
\usepackage[landscape,left=11mm,right=12mm,top=17mm,bottom=17mm]{geometry}
\usepackage{array,xcolor}
\usepackage{amsmath}

\renewcommand{\familydefault}{\sfdefault}
\pagestyle{empty}
\setlength{\parindent}{0pt}
\setlength{\parskip}{0pt}
\setlength{\tabcolsep}{1.6mm}
\setlength{\arrayrulewidth}{0.5pt}
\renewcommand{\arraystretch}{1}
\raggedright

\newcommand{\felt}[1]{{\color{gray}\normalfont\itshape[#1]}}
\ExplSyntaxOn
\NewDocumentCommand{\hent}{m}
  {\prop_item:Nn \g_plan_prop {#1}}
\NewDocumentCommand{\visfelt}{m m}
  {
    \prop_get:NnN \g_plan_prop {#1} \l_tmpa_tl
    \tl_if_blank:VTF \l_tmpa_tl {\felt{#2}} {\tl_use:N \l_tmpa_tl}
  }
\ExplSyntaxOff
\newcolumntype{P}[1]{>{\raggedright\arraybackslash}p{#1}}
\newsavebox{\celleboks}
\newcommand{\mincelle}[2]{%
  \sbox{\celleboks}{\begin{minipage}[t]{\linewidth}\raggedright#2\end{minipage}}%
  \ifdim\dimexpr\ht\celleboks+\dp\celleboks\relax<#1\relax
    \begin{minipage}[t][#1][t]{\linewidth}\raggedright#2\end{minipage}%
  \else\usebox{\celleboks}\fi}
\newcommand{\etikett}[1]{{\bfseries #1}\par}
\newcommand{\sporsmal}[1]{{\itshape #1}\par}
\newlength{\trekolonner}
\setlength{\trekolonner}{\dimexpr\textwidth-6\tabcolsep-4\arrayrulewidth\relax}
\newlength{\sekskolonner}
\setlength{\sekskolonner}{\dimexpr\textwidth-12\tabcolsep-7\arrayrulewidth\relax}
\newlength{\firekolonner}
\setlength{\firekolonner}{\dimexpr\textwidth-8\tabcolsep-5\arrayrulewidth\relax}

\begin{document}

% ==================== SIDE 1 ====================
\vspace*{4mm}
{\centering\fontsize{16}{19}\selectfont\bfseries UNDERVISNINGSPLAN\par}
\vspace{7mm}

\begin{tabular}{|P{.16\firekolonner}|P{.22\firekolonner}|P{.13\firekolonner}|P{.49\firekolonner}|}
\hline
\mincelle{13mm}{\textbf{Dato:} \hent{Dato}} &
\textbf{Trinn:} \hent{Trinn} &
\textbf{Fag:} \hent{Fag} &
\textbf{Varighet:} \hent{Varighet} \\
\hline
\end{tabular}

\vspace{6mm}
\textbf{Kompetansemål:}
\vspace{4mm}
\begin{list}{-}{\setlength{\leftmargin}{12mm}\setlength{\labelsep}{4mm}%
  \setlength{\topsep}{0pt}\setlength{\itemsep}{1pt}\setlength{\parsep}{0pt}}
  \hent{Kompetansemål}
\end{list}
\vspace{19mm}

{\fontsize{9}{10.8}\selectfont
\begin{tabular}{|P{.218\trekolonner}|P{.443\trekolonner}|P{.339\trekolonner}|}
\hline
\mincelle{9mm}{\textbf{Beskrivelse}} &
\textbf{Dine notater} &
\textbf{Begrunnelse:} Hvorfor gjør du dette valget? \\
\hline
\mincelle{35mm}{\etikett{Mål}
\sporsmal{Læringsmål: Operasjonalisering av kompetansemål.}
\vspace{2mm}
\sporsmal{Kunnskaper, ferdigheter og forståelse som elevene skal utvikle, både faglige og sosiale mål.}
\vspace{2mm}
\sporsmal{Faglig innhold.}}
&
\etikett{Læringsmål}
\vspace{2mm}
\visfelt{Læringsmål}{Hva skal elevene kunne, forstå eller mestre etter undervisningen?}
\par\vspace{4mm}
\visfelt{Faglig innhold}{Sentrale begreper og faglig innhold}
&
\visfelt{Begrunnelse for mål}{Begrunn valget av læringsmål og faglig innhold.}
\\
\hline
\mincelle{30mm}{\etikett{Læreforutsetninger/Tilpasning}
\sporsmal{Hvilke forutsetninger har elevene for å lære dette? Hvilke behov for tilpasning ser du?}}
&
\visfelt{Læreforutsetninger}{Forkunnskaper, elevgruppe og behov for tilpasning}
&
\visfelt{Begrunnelse for tilpasning}{Begrunn hvordan undervisningen tilpasses elevene.}
\\
\hline
\end{tabular}
}

\newpage

% ==================== SIDE 2 ====================
{\fontsize{9}{10.8}\selectfont
\begin{tabular}{|P{.218\trekolonner}|P{.443\trekolonner}|P{.339\trekolonner}|}
\hline
\mincelle{29mm}{\etikett{Vurdering og tilbakemelding}
\sporsmal{Hvordan får du informasjon om elevenes læring? Hva skal du se etter? Hva skal eventuelt vurderes og hvordan?}}
&
\visfelt{Vurdering}{Hvordan følger du elevenes læring, og hvordan gir du tilbakemelding?}
&
\visfelt{Begrunnelse for vurdering}{Begrunn vurderingsformene og hva de gir informasjon om.}
\\
\hline
\end{tabular}
}

\vspace{4mm}

% Aktivitetstabellen henter innhold fra del 5 øverst i filen.
{\fontsize{9}{10.8}\selectfont
\begin{tabular}{|P{.185\sekskolonner}|P{.065\sekskolonner}|P{.265\sekskolonner}|P{.12\sekskolonner}|P{.11\sekskolonner}|P{.255\sekskolonner}|}
\hline
\mincelle{17mm}{\etikett{Arbeidsformer og aktiviteter. Læremidler og læringsressurser.}}
& \textbf{Tid}
& \textbf{Hva gjør lærer?}
& \textbf{Hva gjør elevene?}
& \textbf{Læringsressurs:}
& \textbf{Begrunnelse:} \\
\hline
\mincelle{45mm}{\sporsmal{Hvem gjør hva? Når? Hvordan vil du organisere elevenes læringsaktiviteter? Hvilke arbeidsformer vil du benytte? (F.eks. presentasjon, individuelt- eller gruppearbeid, utforskende oppgaver, klassesamtale)}
\vspace{3mm}
\sporsmal{Hvilke læringsressurser vil du bruke? F.eks. lærebok, digitale ressurser)}}
& \visfelt{Tidsplan}{Tid}
& \visfelt{Lærerens aktiviteter}{Lærerens aktiviteter og organisering}
& \visfelt{Elevenes aktiviteter}{Elevenes aktiviteter}
& \visfelt{Læringsressurser}{Materiell og ressurser}
& \visfelt{Begrunnelse for aktiviteter}{Begrunn arbeidsformer, aktiviteter og ressurser.} \\
\hline
\end{tabular}
}

\vspace{8mm}

% Veiledningsfeltet har én ytre ramme som i referansedokumentet.
\begingroup
\setlength{\fboxsep}{1.5mm}
\fbox{\begin{minipage}[t]{\dimexpr\textwidth-2\fboxsep-2\fboxrule\relax}
\mincelle{44mm}{%
\textbf{TIL BRUK I VEILEDNINGSSITUASJONER}
{\fontsize{9}{11}\selectfont (rubrikkene kan utvides etter behov)}
\par\vspace{3mm}
\begin{minipage}[t]{.79\linewidth}\raggedright
Hva ønsker jeg veiledning på?
\par\vspace{5mm}
\visfelt{Spørsmål til veiledning}{Spørsmål og temaer til veiledning}
\end{minipage}\hfill
\begin{minipage}[t]{.20\linewidth}\raggedright
Veileders kommentarer:
\par\vspace{5mm}
\visfelt{Veileders kommentarer}{Kommentarer}
\end{minipage}
}
\end{minipage}}
\endgroup

\end{document}
